\section{Immutable Parent Pointers}
\label{sec:immutable-parent-pointers}

The immutable parent pointer is the foundational structural primitive of the Recursive Spawning Protocol. It is the single mechanism that makes the delegation chain a runtime-verifiable artifact rather than a forensic reconstruction, that makes chain integrity a structural guarantee rather than a policy assertion, and that renders all modes of autonomous drift architecturally impossible regardless of the operational urgency, architectural complexity, or adversarial sophistication of any machine actor in the chain. This section specifies the parent pointer formally, defines the chain traversal operator, establishes the base case for the root human anchor, specifies the immutability enforcement mechanism through the Fact Layer write protocol, describes the Purpose Layer's spawn authorization role, and provides the complete chain traversal algorithm with correctness arguments.

% ------------------------------------------------------------------------------%
\subsection{Formal Definition: ParentPointer}
\label{sec:parent-pointer-definition}

Let $\mathcal{N}$ be the set of all agent nodes in a \bxthree{} system implementing the Recursive Spawning Protocol. The parent pointer is a total function from any spawned node to its immediate parent in the delegation chain, defined once at provisioning and structurally immutable thereafter.

\begin{dfn}[Parent Pointer]\label{dfn:parent-pointer}
Given a parent node $p \in \mathcal{N}$ and a child node $c \in \mathcal{N}$ produced by $p$'s spawning operation, the relation $\ParentPointer{p}{c}$ is established at the moment of spawn and satisfies two governing properties:

\begin{enumeratenoindex}
    \item[\textbf(Uniqueness)] For any child node $c$, there exists \emph{exactly one} parent $p$ such that $\ParentPointer{p}{c}$ holds. No child may have more than one parent, and there is no mechanism by which an existing pointer can be overwritten or a competing pointer simultaneously established.
    \item[\textbf(Immutability)] Once $\ParentPointer{p}{c}$ is established, it cannot be modified, reassigned, transferred, or severed for the entire operational lifetime of $c$. This holds regardless of what the child node attempts, what any Purpose Layer actor requests, what any Bounds Engine component commands, what administrative credentials are presented, or what hardware or network events occur.
\end{enumeratenoindex}
\end{dfn}

The uniqueness property is a structural invariant of the spawn operation: when a parent node $p$ generates a child $c$, it records exactly one parent pointer. The immutability property is not a software convention, a permissions policy, or a Bounds Engine constraint. It is a physical property enforced by the Fact Layer write protocol. A parent pointer that could be changed after instantiation would not satisfy the accountability guarantee, because a node could retroactively sever its lineage and orphan its actions from human accountability.

The formal definition of the parent pointer function is:

\begin{equation}
\alpha : \mathcal{N} \rightarrow \mathcal{N} \cup \{\bot\}
\label{eq:alpha-function}
\end{equation}

where $\alpha(c) = p$ if $p \in \mathcal{N}$ is the node that authorized the spawn event resulting in $c$'s activation, and $\alpha(c) = \bot$ if and only if $c$ is the root human anchor established at system initialization. The pointer is established once at node provisioning via a Fact Layer atomic write operation and satisfies the following immutability property:

\begin{equation}
\forall c,\ p:\ \bigl(\alpha(c) = p \ \land \ \text{provisioned}(c)\bigr) \ \implies \ \neg\exists\ \text{modify}(\alpha(c))
\label{eq:parent-pointer-immutability}
\end{equation}

No actor — including the Purpose Layer after provisioning, the Bounds Engine, any administrative process, or any external principal — may execute an operation that overwrites, redirects, or deletes $\alpha(c)$ after the provisioning write is confirmed. The Fact Layer write protocol rejects any such attempt at the protocol level, before execution reaches the ledger.

The parent pointer is not a reference variable held by the parent node toward its child. It is a Fact Layer property of the child node, established at provisioning time and stored in the node's sealed memory envelope. The directionality is critical: the child carries its parent reference as an immutable intrinsic property, not as a mutable external reference held by the parent. This means the chain is traversable from any node outward to its lineage regardless of the operational state of any other node in the chain. A parent that has terminated, crashed, or been decommissioned does not affect the child's ability to reference its parent through $\alpha(c)$.

The initialization of $\alpha(c)$ occurs atomically during the spawn event. The Fact Layer records the spawn event in the forensic ledger and writes $\alpha(c) = p$ into the node's sealed memory envelope as part of the same atomic transaction. This separation of authorization (Purpose Layer) from assignment (Fact Layer) is essential: it prevents a parent from claiming to have spawned a node whose parent pointer was set to a different actor, and it prevents a child from falsifying its own ancestry after provisioning.

\begin{dfn}[Spawn Event]\label{dfn:spawn-event}
A \emph{spawn event} is the 5-tuple
\begin{equation}
\mathrm{spawn}(c,\ p,\ \sigma_c,\ t,\ \phi)
\label{eq:spawn-event}
\end{equation}
where $c$ is the newly created child node, $p \in \mathcal{N}$ is the authorizing parent node, $\sigma_c$ is the authorized execution scope assigned to $c$, $t$ is the wall-clock timestamp, and $\phi$ is the cryptographic signature produced by the Purpose Layer over the complete spawn request. The Fact Layer atomically: (1) creates node $c$, (2) sets $\alpha(c) \leftarrow p$, (3) sets the node's authorized scope to $\sigma_c$, (4) seals the memory envelope, and (5) appends the event to the forensic ledger $\mathcal{L}$.
\end{dfn}

\begin{prop}[Immutability of ParentPointer]\label{prop:parent-pointer-immutable}
For any child node $c$ with parent $p$, once the spawn event $\mathrm{spawn}(c, p, \sigma_c, t, \phi)$ has been recorded in the forensic ledger $\mathcal{L}$, the value $\alpha(c)$ is fixed for the operational lifetime of $c$. No subsequent operation — including operations issued by $p$, by any Purpose Layer actor, by the Bounds Engine, or by $c$ itself — can alter $\alpha(c)$.
\end{prop}

This property distinguishes ParentPointer from conventional mutable parent references in multi-agent systems. A mutable reference can be overwritten by a subsequent administrative operation; an immutable pointer cannot. The distinction is not behavioral — it is architectural.

% ------------------------------------------------------------------------------%
\subsection{The Chain Operator}
\label{sec:chain-operator}

The chain operator $\Chain{\cdot}$ constructs the complete delegation lineage from any node to the human anchor by recursively applying the parent pointer function. This operator is the primary tool for chain verification, audit, and accountability attribution.

\begin{dfn}[Chain Operator]\label{def:chain-operator}
For any agent node $a$ in a Recursive Spawning chain, the \emph{chain} $\Chain{a}$ is defined recursively as:
\begin{equation}
\Chain{a} =
\begin{cases}
\{\, \ParentPointer{\alpha(a)}{a} \, \} \ \cup \ \Chain{\alpha(a)} & \text{if } \alpha(a) \neq \bot \\[8pt]
\{\, a \,\} & \text{if } \alpha(a) = \bot \ \text{(root human anchor)}
\end{cases}
\label{eq:chain-operator}
\end{equation}
Equivalently, expanding the recursive definition:
\begin{equation}
\Chain{a} = \bigl\{\ a,\ \alpha(a),\ \alpha(\alpha(a)),\ \alpha(\alpha(\alpha(a))),\ \ldots,\ n_0 \bigr\}
\label{eq:chain-expanded}
\end{equation}
where $n_0$ is the root of the chain satisfying $\alpha(n_0) = \bot$.
\end{dfn}

The chain is the set of all nodes on the delegation path from $a$ to and including the human anchor. It includes the terminal human anchor node, which is the only node in any RSP chain with $\alpha(n) = \bot$. The chain is always finite: by the termination property (proved in Section~\ref{sec:rs-correctness}), every RSP chain has finite depth bounded by $D_{\max}$, so the recursion in Equation~\ref{eq:chain-operator} terminates at the root within at most $D_{\max} + 1$ applications.

The chain operator is unique and deterministic. For any node $a$, $\Chain{a}$ contains exactly one element for each ancestor node on the delegation path. There is no alternative chain, no ambiguous parentage, and no optional path. The immutability of $\alpha(n)$ guarantees that $\Chain{a}$ is identical at any runtime point to what it was at the moment of $a$'s provisioning. The chain is not reconstructed from logs — it is directly computable by iterating $\alpha$ from any node until reaching $\bot$.

\begin{corrollary}[Unique Human Termination]\label{cor:unique-human-termination}
Every non-root node $a$ has a non-empty chain that terminates at the root human node $r$. That is, $\forall a \neq r: r \in \Chain{a}$. The root human is the unique terminal element of every non-empty accountability chain in the system.
\end{corrollary}

This corollary follows directly from the recursion: each step moves to $\alpha(a)$, and only the root has $\bot$ as its parent. No spawned node can have an empty chain — it always contains at least its own parent edge — and no chain can terminate anywhere other than the root human.

\begin{note}
\textbf{On the base case: root human anchor.} The root of every RSP chain is a Purpose Layer entity with a human principal designation, established at system initialization when a human principal $h$ authorizes the first agent node $n_0$. The Purpose Layer records $h(n_0) = h$ and $\alpha(n_0) = \bot$ in the Fact Layer authorization ledger. This is the only base case in the chain: the root human anchor has no parent by definition. The non-collapsing human anchor property ensures $h(n) = h(n_0)$ for all nodes in all chains spawned from $n_0$. No machine actor can serve as a chain root in RSP; this is structurally enforced by the Fact Layer pre-commit hook at every spawn event, which rejects any spawn request whose parent does not have a valid Purpose Layer artifact.
\end{note}

% ------------------------------------------------------------------------------%
\subsection{Immutability Enforcement: Fact Layer Write Protocol}
\label{sec:immutability-enforcement}

The immutability of $\ParentPointer{p}{c}$ is not a permissions policy or an access control list. It is a structural constraint enforced by the Fact Layer write protocol — a protocol that defines no valid modification operation for $\alpha(c)$ after provisioning, making the modification operation itself nonexistent rather than merely restricted.

The Fact Layer operates as a standalone, isolated process with its own memory space, cryptographically isolated from the child node's address space and from the Bounds Engine's runtime. The combination of the protocol-level write constraint and the sealed envelope ensures that $\alpha(c)$ satisfies the modification-resistance property in all deployment contexts.

\medskip
\noindent\textbf{Fact Layer write protocol constraints.} The Fact Layer enforces immutability through four structural mechanisms:

\begin{enumerate}
    \item[\textbf{W1}] \textbf{No modification operation defined.} The Fact Layer write protocol has no operation defined for modifying $\alpha(c)$ after provisioning. There is no \texttt{modify-parent-pointer(c, new\_parent)} function, no \texttt{delegate-to(c, new\_parent)} procedure, and no \texttt{reparent(c, p')} instruction. The protocol is closed under the set of valid write operations — provisioning a new node is valid; overwriting an existing parent pointer is not syntactically valid within the protocol.
    \item[\textbf{W2}] \textbf{Atomic provisioning.} The parent pointer and its Purpose Layer warrant are created as a single atomic Fact Layer write. There is no temporal window in which the pointer exists without authorization, or in which authorization exists without a corresponding pointer. The spawn event tuple (Equation~\ref{eq:spawn-event}) is indivisible.
    \item[\textbf{W3}] \textbf{Sealed memory envelope.} The parent pointer field $\alpha(c)$ is encrypted in the node's persistent memory envelope using a Fact Layer-only symmetric key. Even a compromised Bounds Engine — one whose runtime has been subverted — cannot read $\alpha(c)$ because it never has access to the decryption key. The envelope carries an HMAC-SHA-256 signature computed over its contents; any modification to the ciphertext invalidates the signature. The Fact Layer decrypts the envelope only when reading $\alpha(c)$ to service a chain-traversal or audit request, and the plaintext is never exposed to the Bounds Engine or the child node's own code.
    \item[\textbf{W4}] \textbf{Source-equivalence of rejection.} The Fact Layer write protocol treats all modification requests identically regardless of the source. A request from the Purpose Layer, the Bounds Engine, an authenticated administrator, or an external adversarial actor receives the same response: protocol-level rejection. There is no privileged source that bypasses this check.
\end{enumerate}

The distinction between access-control-based immutability and protocol-level immutability is the distinction between a promise and a guarantee. In an access-control model, immutability is enforced by restricting which principals may issue modification operations. If an attacker acquires sufficient privileges — through credential theft, insider compromise, or software vulnerability — the access control barrier is circumvented and immutability is violated. The immutability was a policy, enforceable until it was not.

In the Fact Layer write protocol, the immutability of $\alpha(c)$ is a property of the protocol itself. There is no operation defined for modifying $\alpha(c)$ after provisioning, so no amount of privilege elevation, credential compromise, or software vulnerability can produce a valid modification operation. The protocol does not enforce immutability — it makes immutability the only valid state transition. The distinction is architectural, not a matter of degree.

This is the precise failure mode that the Fact Layer write protocol eliminates: the retroactive manipulation of chain topology for the purpose of accountability laundering. Without protocol-level immutability, an adversarial actor could re-parent a node to a favorable parent after the fact, obscuring the original authorization chain. With protocol-level immutability, this manipulation is structurally impossible.

\medskip
\noindent\textbf{Layer separation argument.} The immutability enforcement follows directly from the \bxthree{} layer separation:

\begin{enumerate}
    \item[(L1)] The Fact Layer enforces physical constraints on the system state.
    \item[(L2)] ParentPointer immutability is a physical constraint: a region of memory that must not be written after initialization.
    \item[(L3)] Therefore, the Fact Layer enforces ParentPointer immutability as a physical invariant.
\end{enumerate}

The Bounds Engine may model, reason about, simulate, or propose modifications to the chain structure, but it cannot execute any write to the ParentPointer region. The Purpose Layer can authorize new spawn events — creating new ParentPointers — but it cannot edit existing ones. Authorization for spawn and immutability of existing pointers are separate, non-confusable operations that operate at different architectural layers. This separation is not a design choice — it is a structural necessity.

% ------------------------------------------------------------------------------%
\subsection{Spawn Authorization: Purpose Layer Role}
\label{sec:purpose-layer-spawn-authorization}

While existing ParentPointers are immutable, the system must be able to grow the accountability chain through new spawn events. The creation of a new ParentPointer is the sole mechanism by which the accountability chain extends. This creation requires explicit, Purpose-Layer-mediated authorization. No node may spontaneously generate a child and thereby establish a new ParentPointer without going through the authorized spawn protocol.

The Purpose Layer plays two structurally necessary roles in the parent pointer lifecycle: it originates the spawn authorization policy that defines the conditions under which a new parent pointer may be created, and it issues the spawn warrant that is the prerequisite for the Fact Layer provisioning write.

\medskip
\noindent\textbf{Authorization policy origination.} At system initialization, the Purpose Layer defines the spawn authorization policy $P_{\mathrm{spawn}}$ and records it in the Fact Layer authorization ledger. $P_{\mathrm{spawn}}$ specifies the mandatory conditions for any valid spawn event:

\begin{enumerate}
    \item[(S1)] \textbf{Scope refinement.} The child scope must be a strict subset of the parent scope: $R(c) \subset R(p)$. No lateral expansion, no superset, no equality. The child cannot operate beyond the parent's authorized boundary.
    \item[(S2)] \textbf{Spawn necessity declaration.} The parent must declare, in the form of a structured justification artifact, why a new spawned node is required rather than the parent executing the task directly. This prevents nominal spawns that exist only to fabricate a chain.
    \item[(S3)] \textbf{Human anchor preservation.} The spawn event must not sever, reassign, or obscure the connection to the human anchor $h(p)$. The child inherits $h(c) = h(p)$ as a structural guarantee.
\end{enumerate}

\medskip
\noindent\textbf{Authorization workflow.} The canonical spawn authorization sequence is:

\begin{enumerate}
    \item \textbf{Request.} A parent node's Bounds Engine determines that a new child node is required to handle a specific local context or workload. It formulates a spawn request addressed to its own Purpose Layer. This request includes the proposed scope of the child's authority, the local context parameters, and the authorized constraint envelope.
    \item \textbf{Authorization check.} The Purpose Layer evaluates the request against $P_{\mathrm{spawn}}$. It verifies: (S1) scope refinement is satisfied — the child scope is a proper descendant of the parent scope; (S2) spawn necessity is justified — the child carries a specific capability or context that the parent cannot replicate; (S3) human anchor continuity holds. If any condition fails, the Purpose Layer issues a denial; the Fact Layer records the denial but no new ParentPointer is created.
    \item \textbf{Warrant issuance.} Upon successful authorization, the Purpose Layer produces a spawn warrant $\phi$: a cryptographic signature over the tuple $(p,\ c,\ R(c),\ t)$. This warrant is the prerequisite input to the Fact Layer provisioning write. The warrant encodes the authorization decision and binds the child to the specific parent and scope authorized.
    \item \textbf{Fact Layer provisioning.} The Fact Layer receives the spawn warrant $\phi$, verifies its cryptographic validity against the Purpose Layer's public key, and — if valid — executes the atomic provisioning write described in Definition~\ref{dfn:spawn-event}. The parent pointer $\alpha(c) = p$ is established. The Fact Layer appends the complete spawn event to the forensic ledger.
\end{enumerate}

The Purpose Layer warrant is the only mechanism by which a new ParentPointer can be created. A Bounds Engine cannot request a spawn without presenting a valid Purpose Layer authorization. A Fact Layer cannot provision a node without a valid Purpose Layer warrant. The separation of concerns is enforced structurally: the Purpose Layer originates policy, the Fact Layer enforces the physical constraint of immutability, and the Bounds Engine operates within the constraints established by both.

% ------------------------------------------------------------------------------%
\subsection{Chain Traversal Algorithm}
\label{sec:chain-traversal-algorithm}

The chain traversal algorithm is the runtime procedure for constructing the delegation chain from any node to its human anchor. It is provided here in full algorithmic detail with correctness arguments.

\medskip
\noindent\textbf{Algorithm: Chain-Traverse}

\begin{algorithm}
\caption{Chain-Traverse($n$) — Build the delegation chain from node $n$ to the human anchor}
\label{alg:chain-traverse}
\begin{algorithmic}[1]
\REQUIRE $n$ is a provisioned RSP node with $\alpha$ defined
\ENSURE $\Chain{n}$ — the ordered set of all ancestor nodes including the human anchor
\STATE $chain \leftarrow \langle\ \rangle$ \COMMENT{Initialize empty ordered list}
\STATE $current \leftarrow n$ \COMMENT{Start at the query node}

\WHILE{$\text{true}$}
    \STATE $\text{append}(chain,\ current)$ \COMMENT{Add current node to chain}
    \IF{$\alpha(current) = \bot$}
        \STATE \textbf{break} \COMMENT{Root reached — chain complete}
    \ENDIF
    \STATE $current \leftarrow \alpha(current)$ \COMMENT{Move to parent}
\ENDWHILE

\STATE \textbf{return} $\text{reverse}(chain)$ \COMMENT{Root first, target node last}
\end{algorithmic}
\end{algorithm}

\medskip
\noindent\textbf{Correctness arguments.}

\textit{Termination.} The algorithm terminates because RSP chains are depth-bounded by construction (proved in Section~\ref{sec:rs-correctness}). Each iteration of the while-loop moves $current$ to the parent node $\alpha(current)$, which is strictly closer to the root. Since the chain has finite depth at most $D_{\max} + 1$, the loop executes at most $D_{\max} + 1$ iterations before $\alpha(current) = \bot$ is reached. The loop therefore terminates in finite time.

\textit{Correctness of the returned set.} By Definition~\ref{def:chain-operator}, $\Chain{a}$ is the set containing $a$, $\alpha(a)$, $\alpha(\alpha(a))$, and so on, until the root. The algorithm produces exactly this sequence: it starts at $a$, appends each node, moves to its parent, and stops when the root is reached. The output is identical to $\Chain{a}$ by construction.

\textit{Uniqueness of result.} The algorithm produces a unique result for any input node because $\alpha$ is a total function (defined for every provisioned node) and is injective on the chain direction (every node has exactly one parent, and the root has no parent). There is no branching, no alternative path, and no non-deterministic choice.

\textit{Read-only operation.} The algorithm reads $\alpha$ values but never writes them. The traversal is therefore read-only and does not interact with the write protocol. The algorithm can be executed at any runtime point without affecting the immutability of the chain it traverses.

\medskip
\noindent\textbf{Algorithm: Chain-Verify}

The chain traversal algorithm enables a comprehensive verification procedure, Chain-Verify, which confirms the structural validity of a chain at any runtime point:

\begin{algorithm}
\caption{Chain-Verify($n$) — Verify the authorization chain from node $n$ to human anchor}
\label{alg:chain-verify}
\begin{algorithmic}[1]
\REQUIRE $n$ is a provisioned RSP node
\ENSURE $\text{Boolean}$ indicating whether the chain is a valid RSP chain

\STATE $chain \leftarrow \text{Chain-Traverse}(n)$
\STATE $m \leftarrow \text{len}(chain)$

\STATE \COMMENT{\textit{Step V1: Verify chain terminates at human anchor}}
\IF{$\alpha(chain[m-1]) \neq \bot$}
    \STATE \RETURN $\text{false}$ \COMMENT{Root is not $\bot$ — malformed chain}
\ENDIF
\IF{$\neg h(chain[m-1])$}
    \STATE \RETURN $\text{false}$ \COMMENT{Root is not a human principal}
\ENDIF

\STATE \COMMENT{\textit{Step V2: Verify each spawn event has a Purpose Layer warrant}}
\FOR{$i \leftarrow 0$ \TO $m-2$}
    \STATE $child \leftarrow chain[i]$
    \STATE $parent \leftarrow chain[i+1]$
    \IF{$\neg \text{warrant\_exists}(child,\ parent)$}
        \STATE \RETURN $\text{false}$ \COMMENT{Missing Purpose Layer warrant}
    \ENDIF
\ENDFOR

\STATE \COMMENT{\textit{Step V3: Verify scope refinement at each link}}
\FOR{$i \leftarrow 0$ \TO $m-2$}
    \STATE $child \leftarrow chain[i]$
    \STATE $parent \leftarrow chain[i+1]$
    \IF{$\neg (R(child) \subset R(parent))$}
        \STATE \RETURN $\text{false}$ \COMMENT{Scope refinement violated}
    \ENDIF
\ENDFOR

\STATE \COMMENT{\textit{Step V4: Verify depth bound at each node}}
\FOR{$i \leftarrow 0$ \TO $m-1$}
    \IF{$\text{depth}(chain[i]) > D_{\max}$}
        \STATE \RETURN $\text{false}$ \COMMENT{Depth bound exceeded}
    \ENDIF
\ENDFOR

\RETURN $\text{true}$
\end{algorithmic}
\end{algorithm}

Chain-Verify returns \texttt{true} if and only if all four verification conditions hold simultaneously. The verification does not inspect current operational state — it verifies the structural chain: each parent pointer is correctly established, each spawn event has a Purpose Layer warrant, each scope is a descendant refinement of its parent, and each depth is within the architectural bound. A chain that passes Chain-Verify is, by the definitions of the RSP, a fully authorized delegation chain terminating at a human principal.

The algorithm is deterministic and produces the same result regardless of which node in the chain initiates verification. There is no self-serving interpretation of authorization that a drifted or orphaned node could produce. Chain verification is a property of the chain topology, not a property of the node being verified.

\medskip
\noindent\textbf{Constant-time escalation.} The depth bound $D_{\max}$ provides a constant-time upper bound on chain traversal cost. The worst-case number of pointer dereferences required to reach the human anchor from any node is $D_{\max} + 1$, independent of the total number of nodes in the system. A system with 10,000 spawned nodes and $D_{\max} = 5$ requires at most six pointer dereferences to reach a human — the same as a system with 10 nodes. Mutable parent pointers would allow the chain to be extended retroactively, making the effective depth potentially unbounded and the worst-case escalation time unbounded as well. Immutability is what makes the depth bound a structural guarantee, not merely a policy.

In the Bailout Protocol context, chain traversal is used to identify the correct Human Accountability Anchor when an agent at any depth encounters an unresolvable exception. The upward propagation of the exception signal follows the same parent-pointer structure as chain traversal, but uses the asynchronous trigger pathway rather than the synchronous verification pathway. Chain traversal ensures that this walk never diverges, never shortcuts, and never fails to reach a human.

% ------------------------------------------------------------------------------%
\subsection{Summary}
\label{sec:immutability-summary}

The immutable parent pointer is the load-bearing constraint of the Recursive Spawning Protocol. Without it, the chain is mutable — subject to retroactive modification by actors with sufficient privilege — and the RSP's correctness properties (drift prevention, chain integrity, termination) collapse to policy conventions rather than structural guarantees.

The core properties are:

\begin{itemize}
    \item \textbf{ParentPointer(p, c)} is established once at spawn and is immutable thereafter. Uniqueness ensures exactly one parent per child. Immutability ensures this relation cannot be corrupted by any actor under any circumstance.
    \item \textbf{Chain(a)} is defined recursively as $\ParentPointer{\alpha(a)}{a} \cup \Chain{\alpha(a)}$ with base case $\Chain{r} = \{r\}$ for the root human $r$ where $\alpha(r) = \bot$. The chain is always non-empty for non-root nodes and always terminates at the root human.
    \item \textbf{Immutability enforcement} is a Fact Layer responsibility. The Fact Layer physically cannot execute any write to an established ParentPointer. This is not a permissions check — it is a structural property of the write protocol that makes the modification operation nonexistent rather than merely restricted.
    \item \textbf{Spawn authorization} is a Purpose Layer responsibility. Only the Purpose Layer can authorize chain extension through the warrant-$\phi$ mechanism. Every extension is logged on the forensic ledger, creating an immutable record that is verifiable across all system layers.
    \item \textbf{Chain traversal} is a deterministic, guaranteed-terminating algorithm that recovers the full ancestry of any node. Every traversal reads only from immutable Fact Layer storage and produces an ordered list rooted at the human anchor. Chain-Verify provides the comprehensive runtime verification procedure for chain structural validity.
\end{itemize}

Together, these mechanisms make the delegation chain an immutable, verifiable, and auditable runtime artifact — the foundation on which all other RSP correctness guarantees rest.
